
This paper investigates whether community breadth at benchmark introduction --- operationalized as the count of distinct paper submissions per benchmark through its plurality-introduction year on Papers With Code --- predicts plurality benchmark displacement hazard. Two competing mechanisms motivate the hypothesis: a lock-in mechanism (broad early adoption creates investment network effects that resist displacement, HR < 1) and a saturation mechanism (widespread adoption signals a solved benchmark, accelerating replacement, HR > 1). Before testing the mechanism, a 5-gate FAIL FAST pre-validation protocol confirms data quality: join coverage (G0 = 0.862), predictor time-independence via partial $R^2$ (G1 = 0.605, G2 = 0.975), adequate cross-benchmark variance (G3 std = 0.246), and absence of multicollinearity (G4 max VIF = 2.14). All five gates pass by wide margins, establishing that any downstream null reflects genuine absence of effect rather than data quality failure. Cox proportional hazards regression (CoxPHFitter, L2 penalizer = 0.1) on 258 complete-case rows from the h-e2 panel (87 Koch et al. 2021 tasks, 2015--2023; 176 displacement events) yields HR = 1.006 (95\% CI = [0.846, 1.196]), LRT p = 0.9495 --- a near-perfect null. The diversity predictor adds $\Delta logL$ = 0.0020 to the baseline model. Neither the lock-in nor the saturation mechanism is detectable at this operationalization. Submission-count community breadth diversity does not predict benchmark displacement timing in the Papers With Code ecosystem. The 5-gate FAIL FAST protocol constitutes a reusable methodological contribution for benchmark lifecycle survival analysis; score trajectory and institutional diversity remain open directions for future investigation.

